%% Drafted from the mapped corpus by scripts/draft_topics.py.
% Model: gpt-5.6-terra; source characters: 269101.
\section{Основни понятия}

Методът на резолюциите е доказателствен метод, който свежда проверката за логическо следване до проверка за неизпълнимост. Идеята е проста: за да се докаже, че от множество формули $\Gamma$ следва формулата $\varphi$, разглеждаме множеството
\[
\Gamma\cup\{\neg\varphi\}.
\]
Ако то е неизпълнимо, то няма структура, в която всички формули от $\Gamma$ са верни, а $\varphi$ е невярна; следователно $\Gamma\models\varphi$. Резолюцията търси синтактично свидетелство за тази неизпълнимост: извеждане на празния дизюнкт.

\begin{defn}
Литерал е атомарна формула или отрицание на атомарна формула. Ако $L$ е литерал, то противоположният му литерал ще означаваме с $\overline L$:
\[
\overline{P(t_1,\ldots,t_n)}=\neg P(t_1,\ldots,t_n),
\qquad
\overline{\neg P(t_1,\ldots,t_n)}=P(t_1,\ldots,t_n).
\]
\end{defn}

\begin{defn}
Елементарна дизюнкция е безкванторна формула, изградена от литерали само с дизюнкция. Формула е в \emph{конюнктивна нормална форма} (КНФ), ако е конюнкция от елементарни дизюнкции.
\end{defn}

В резолютивния метод не е необходимо да се пазят редът и повторенията на литералите в една дизюнкция. Затова дизюнкциите се представят като множества.

\begin{defn}
Дизюнкт е крайно множество от литерали. Допуска се и празният дизюнкт, означаван с $\Box$ или $\varnothing$.
\end{defn}

Например формулата
\[
P(x)\vee\neg Q(f(x))\vee P(x)
\]
съответства на дизюнкта
\[
\{P(x),\neg Q(f(x))\}.
\]
Празният дизюнкт няма литерали и съответства на лъжа: той е неверен при всяка оценка.

\begin{defn}
Нека $\mathcal M$ е структура, а $v$ е оценка на индивидуалните променливи. Дизюнктът $D$ е \emph{верен в $\mathcal M$ при $v$}, ако поне един негов литерал е верен:
\[
\mathcal M\models_v D
\quad\Longleftrightarrow\quad
(\exists L\in D)\ \mathcal M\models_v L.
\]
Той е неверен в $\mathcal M$ при $v$, ако всички негови литерали са неверни.
\end{defn}

\begin{defn}
Дизюнктът $D$ е \emph{тъждествено верен в структурата} $\mathcal M$, ако
\[
(\forall v)\ \mathcal M\models_v D.
\]
Множество от дизюнкти $\Gamma$ е \emph{изпълнимо}, ако съществува структура $\mathcal M$, в която всеки дизюнкт от $\Gamma$ е тъждествено верен:
\[
\mathcal M\models\Gamma
\quad\Longleftrightarrow\quad
(\forall D\in\Gamma)\,(\forall v)\ \mathcal M\models_v D.
\]
Тогава $\mathcal M$ се нарича модел на $\Gamma$.
\end{defn}

Важно е да се различават дизюнктът $\{p,\neg p\}$ и множеството от формули $\{p,\neg p\}$. Първият е тавтологичен дизюнкт, понеже съдържа противоположни литерали. Второто множество е неизпълнимо, понеже изисква едновременно $p$ и $\neg p$ да са верни.

\begin{remark}
Ако формула е приведена в КНФ
\[
C_1\wedge\cdots\wedge C_m,
\]
то на нея съответства множеството от дизюнкти $\{D_1,\ldots,D_m\}$, където $D_i$ е множеството от литералите в $C_i$. Изпълнимостта на формулата е еквивалентна на изпълнимостта на съответното множество от дизюнкти.
\end{remark}

\section{Резолюция за дизюнкти}

\subsection{Същинска резолвента}

Най-напред разглеждаме случая, в който два дизюнкта съдържат буквално противоположни литерали.

\begin{defn}
Нека $D_1$ и $D_2$ са дизюнкти, $L\in D_1$ и $\overline L\in D_2$. Дизюнктът
\[
R_L(D_1,D_2)
=
(D_1\setminus\{L\})\cup(D_2\setminus\{\overline L\})
\]
се нарича \emph{непосредствена резолвента} на $D_1$ и $D_2$ относно литерала $L$.
\end{defn}

Например:
\[
\{p,q\},\qquad \{\neg p,r\}
\]
имат резолвента
\[
\{q,r\}.
\]
При резолюция се премахва точно една двойка противоположни литерали. Ако два дизюнкта съдържат повече от една противоположна двойка, не се премахват всички двойки едновременно.

\begin{example}
За дизюнктите
\[
D_1=\{p,q,r\},
\qquad
D_2=\{\neg p,\neg q,r\}
\]
резолвентите са
\[
\{q,\neg q,r\}
\quad\text{и}\quad
\{p,\neg p,r\},
\]
но не и $\{r\}$.
\end{example}

\begin{lem}[Резолютивна лема]
Нека $D$ е резолвента на $D_1$ и $D_2$. За всяка булева интерпретация $I$ е вярно:
\[
I\models D_1\ \text{и}\ I\models D_2
\quad\Longrightarrow\quad
I\models D.
\]
\end{lem}

\begin{proof}
Нека
\[
D=(D_1\setminus\{L\})\cup(D_2\setminus\{\overline L\}).
\]
Разглеждаме два случая.

Ако $I\models L$, то $I\not\models\overline L$. Понеже $I\models D_2$, някой литерал от $D_2$, различен от $\overline L$, трябва да е верен. Той принадлежи на $D$, следователно $I\models D$.

Ако $I\not\models L$, то понеже $I\models D_1$, някой литерал от $D_1$, различен от $L$, е верен. Той принадлежи на $D$, следователно $I\models D$.
\end{proof}

\subsection{Резолютивен извод}

\begin{defn}
Нека $\Gamma$ е множество от дизюнкти. \emph{Резолютивен извод} на дизюнкта $\delta$ от $\Gamma$ е крайна редица
\[
\delta_0,\delta_1,\ldots,\delta_n,
\qquad \delta_n=\delta,
\]
такава че за всеки $\delta_i$ е изпълнено поне едно от условията:

\begin{enumerate}
\item $\delta_i\in\Gamma$;
\item $\delta_i$ е резолвента на два дизюнкта, срещащи се по-рано в редицата.
\end{enumerate}

Казваме, че $\delta$ се извежда от $\Gamma$ с резолюция и пишем
\[
\Gamma\vdash_r\delta.
\]
\end{defn}

\begin{keythm}[Коректност на резолютивната изводимост]
Ако
\[
\Gamma\vdash_r\Box,
\]
то $\Gamma$ е неизпълнимо.
\end{keythm}

\begin{proof}
Нека
\[
\delta_0,\delta_1,\ldots,\delta_n=\Box
\]
е резолютивен извод от $\Gamma$. Ще докажем по индукция по $i$, че всеки модел на $\Gamma$ удовлетворява $\delta_i$.

Нека $\mathcal M\models\Gamma$. Ако $\delta_i\in\Gamma$, то по дефиниция
\[
\mathcal M\models\delta_i.
\]
Ако $\delta_i$ е резолвента на два по-ранни дизюнкта $\delta_j$ и $\delta_k$, то по индукционната хипотеза
\[
\mathcal M\models\delta_j
\quad\text{и}\quad
\mathcal M\models\delta_k.
\]
От резолютивната лема следва
\[
\mathcal M\models\delta_i.
\]

Следователно всеки модел на $\Gamma$ би трябвало да удовлетворява и $\Box$. Но празният дизюнкт няма верен литерал и е неверен при всяка интерпретация. Следователно $\Gamma$ няма модел.
\end{proof}

\begin{example}
Нека
\[
\Gamma=
\bigl\{
\{p,q\},
\{\neg p,q\},
\{p,\neg q\},
\{\neg p,\neg q\}
\bigr\}.
\]
Получаваме извода
\[
\begin{array}{rll}
1.&\{p,\neg q\} & \in\Gamma,\\
2.&\{p,q\} & \in\Gamma,\\
3.&\{p\} & \text{резолвента на 1 и 2 относно }q,\\
4.&\{\neg p,\neg q\} & \in\Gamma,\\
5.&\{\neg p,q\} & \in\Gamma,\\
6.&\{\neg p\} & \text{резолвента на 4 и 5 относно }q,\\
7.&\Box & \text{резолвента на 3 и 6 относно }p.
\end{array}
\]
Следователно $\Gamma$ е неизпълнимо.
\end{example}

\section{Предикатна либерална резолюция}

При предикатните дизюнкти противоположните литерали невинаги са синтактично еднакви. Например $P(x)$ и $\neg P(f(y))$ могат да станат противоположни след подходяща субституция. Затова е необходима унификация.

\subsection{Субституции и унификация}

\begin{defn}
Субституция е крайно множество
\[
\sigma=
\left\{
\frac{x_1}{t_1},\ldots,\frac{x_n}{t_n}
\right\},
\]
където $x_1,\ldots,x_n$ са различни индивидуални променливи, а $t_1,\ldots,t_n$ са термове. Резултатът от прилагането на $\sigma$ върху дизюнкта $\varepsilon$ се означава с $\varepsilon\sigma$.
\end{defn}

\begin{defn}
Субституцията $\sigma$ е унификатор на термовете $t_1,\ldots,t_m$, ако
\[
t_1\sigma=\cdots=t_m\sigma.
\]
Тя е най-общ унификатор, ако всеки друг унификатор се получава от нея чрез последваща субституция.
\end{defn}

При унификацията трябва да се спазва проверката за участие: не може да се замества $x$ с терм, в който участва самото $x$. Например уравнението
\[
x=f(x)
\]
няма допустим унификатор.

\begin{example}
Литералите
\[
P(x,f(y))
\quad\text{и}\quad
\neg P(f(z),f(a))
\]
се унифицират чрез
\[
\sigma=
\left\{
\frac{x}{f(z)},
\frac{y}{a}
\right\}.
\]
След прилагане на $\sigma$ получаваме противоположните литерали
\[
P(f(z),f(a))
\quad\text{и}\quad
\neg P(f(z),f(a)).
\]
\end{example}

\subsection{Предикатна либерална резолвента}

\begin{defn}
Нека $\varepsilon_1$ и $\varepsilon_2$ са дизюнкти. Първо техните променливи се преименуват така, че двата дизюнкта да нямат общи променливи. Нека
\[
P(t_1,\ldots,t_n)\in\varepsilon_1,
\qquad
\neg P(s_1,\ldots,s_n)\in\varepsilon_2,
\]
а $\sigma$ е най-общ унификатор на съответните аргументи. Тогава дизюнктът
\[
\left(
(\varepsilon_1\setminus\{P(t_1,\ldots,t_n)\})
\cup
(\varepsilon_2\setminus\{\neg P(s_1,\ldots,s_n)\})
\right)\sigma
\]
се нарича \emph{предикатна либерална резолвента} на $\varepsilon_1$ и $\varepsilon_2$.
\end{defn}

Терминът ``либерална'' подчертава, че преди резолюцията литералите могат да се направят противоположни чрез унификация, а не е необходимо да са буквално противоположни още в изходните дизюнкти.

\begin{example}
Нека
\[
\varepsilon_1=\{P(x),Q(x)\},
\qquad
\varepsilon_2=\{\neg P(f(y)),R(y)\}.
\]
За $P(x)$ и $\neg P(f(y))$ най-общ унификатор е
\[
\sigma=\left\{\frac{x}{f(y)}\right\}.
\]
Следователно либералната резолвента е
\[
\{Q(f(y)),R(y)\}.
\]
\end{example}

\begin{defn}
Нека $\Gamma$ е множество от предикатни дизюнкти. \emph{Предикатен резолютивен извод} на $\delta$ от $\Gamma$ е крайна редица от дизюнкти, в която всеки член е или елемент на $\Gamma$, или предикатна либерална резолвента на два по-ранни члена. Отново пишем
\[
\Gamma\vdash_r\delta.
\]
\end{defn}

\begin{supp}[Бележка]
В някои представяния предикатната резолюция се описва чрез дизюнкти с ограничения, които записват равенствата между аргументите на резолвираните литерали. След елиминиране на ограниченията чрез алгоритъм за унификация се получава същата резолвента с приложена субституция. Тук използваме непосредственото описание чрез най-общ унификатор.
\end{supp}

\subsection{Коректност и пълнота}

\begin{keythm}[Коректност на предикатната резолютивна изводимост]
Ако
\[
\Gamma\vdash_r\Box,
\]
то множеството от предикатни дизюнкти $\Gamma$ е неизпълнимо.
\end{keythm}

Доказателствената идея е същата като при същинската резолюция. Ако дадена структура удовлетворява двата родителски дизюнкта при всички оценки, то след прилагане на унифициращата субституция тя удовлетворява и резолвентата. Следователно всеки дизюнкт в резолютивния извод е логическо следствие на изходното множество. Появата на $\Box$ противоречи на съществуването на модел.

\begin{keythm}[Пълнота на резолютивната изводимост]
За множество $\Gamma$ от първоредови дизюнкти в език без формално равенство е в сила:
\[
\Gamma\ \text{е неизпълнимо}
\quad\Longleftrightarrow\quad
\Gamma\vdash_r\Box.
\]
\end{keythm}

При същинската резолюция пълнотата може да се изрази чрез резолютивното затваряне. Нека
\[
R(S)=S\cup
\{D\mid D\text{ е резолвента на два дизюнкта от }S\},
\]
\[
S_0=S,\qquad S_{n+1}=R(S_n),\qquad
S^*=\bigcup_{n\geq 0}S_n.
\]
Тогава
\[
S\vdash_rD
\quad\Longleftrightarrow\quad
D\in S^*,
\]
а в частност
\[
S\text{ е неизпълнимо}
\quad\Longleftrightarrow\quad
\Box\in S^*
\quad\Longleftrightarrow\quad
S\vdash_r\Box.
\]

За предикатния случай пълнотата е свързана с теоремата на Ербран. След сколемизация и преминаване към затворени универсални формули се разглеждат техните затворени частни случаи. Неизпълнимостта на първоредовото множество се свежда до булева неизпълнимост на подходящо крайно множество от такива случаи, за което резолюцията е пълна.

Ако равенството е част от логическия език със стандартната си семантика, обикновената резолюция сама не е достатъчна за тази формулировка на пълнотата. Необходими са подходящи аксиоми за равенство или специализирани правила като параметодулация/суперпозиция.

\begin{remark}
Пълнотата не означава, че всяка стратегия за избор на двойки дизюнкти непременно ще намери празния дизюнкт за краен брой стъпки. Теоремата твърди, че при неизпълнимо множество съществува краен резолютивен извод на $\Box$.
\end{remark}

\section{Компютърно генериране на доказателства}

Резолюцията е особено подходяща за автоматизация, защото всяка стъпка има краен и проверим вид: посочват се два родителски дизюнкта, избрани литерали и унификатор. Така доказателството може да се пази като дърво или като линейна последователност от стъпки.

Общият процес има следните етапи:

\begin{enumerate}
\item Формулира се задачата като проверка на неизпълнимост, обикновено чрез добавяне на отрицанието на целта.
\item Отстраняват се импликациите и еквивалентностите.
\item Отрицанията се свеждат непосредствено пред атомарни формули.
\item Формулата се привежда в пренексна форма и се сколемизира, когато има квантори за съществуване.
\item Отпадат универсалните квантори и матрицата се привежда в КНФ.
\item Генерират се резолвенти чрез унификация.
\item При получаване на $\Box$ се извежда неизпълнимостта на изходното множество.
\end{enumerate}

\begin{example}[Предикатно доказване на неизпълнимост]
Нека са дадени формулите
\[
\forall x\,\bigl(Student(x)\Rightarrow Learns(x)\bigr),
\]
\[
Student(ivan),
\]
\[
\neg Learns(ivan).
\]
След преминаване към дизюнкти получаваме:
\[
D_1=\{\neg Student(x),Learns(x)\},
\qquad
D_2=\{Student(ivan)\},
\qquad
D_3=\{\neg Learns(ivan)\}.
\]
Резолюцията на $D_1$ и $D_2$ използва унификатора
\[
\sigma=\left\{\frac{x}{ivan}\right\}
\]
и дава
\[
D_4=\{Learns(ivan)\}.
\]
След това $D_4$ и $D_3$ дават
\[
\Box.
\]
Следователно трите изходни формули са неизпълними като множество.
\end{example}

\section{Хорнови дизюнкти и Пролог}

\begin{defn}
Хорнов дизюнкт е дизюнкт с най-много един положителен литерал.
\end{defn}

Най-често се разглеждат следните три вида:

\[
\{P\}
\]
е факт;

\[
\{P,\neg Q_1,\ldots,\neg Q_n\}
\]
е правило, което се чете като
\[
Q_1\wedge\cdots\wedge Q_n\Rightarrow P;
\]

\[
\{\neg Q_1,\ldots,\neg Q_n\}
\]
е цел, която съответства на въпроса дали могат да се докажат едновременно $Q_1,\ldots,Q_n$.

Хорновите правила и факти образуват хорнова програма. Резолюцията за хорнови дизюнкти е основата на логическото програмиране и на езика Пролог. Вместо да се генерират произволни резолвенти, обикновено се работи с цел и с правило, чиято глава се унифицира с избран подцелеви атом.

\begin{example}[Дефиниране на предикати на Пролог]
Следната програма описва предиката \verb|predok(X,Y)|: ``\verb|X| е предшественик на \verb|Y|''.

\begin{verbatim}
roditel(ivan, maria).
roditel(maria, petar).

predok(X, Y) :- roditel(X, Y).
predok(X, Y) :- roditel(X, Z), predok(Z, Y).
\end{verbatim}

Фактът \verb|roditel(ivan,maria).| съответства на едноелементния дизюнкт
\[
\{roditel(ivan,maria)\}.
\]
Правилото
\begin{verbatim}
predok(X, Y) :- roditel(X, Z), predok(Z, Y).
\end{verbatim}
съответства на хорновия дизюнкт
\[
\{
predok(X,Y),
\neg roditel(X,Z),
\neg predok(Z,Y)
\}.
\]
\end{example}

\begin{example}[Резолютивно доказване на цел]
Нека целта е
\begin{verbatim}
?- predok(ivan, petar).
\end{verbatim}
Тя се представя с отрицателния дизюнкт
\[
\{\neg predok(ivan,petar)\}.
\]
Резолюция с рекурсивното правило дава нова цел
\[
\{
\neg roditel(ivan,Z),
\neg predok(Z,petar)
\}.
\]
Фактът $roditel(ivan,maria)$ унифицира $Z$ с $maria$, откъдето получаваме
\[
\{\neg predok(maria,petar)\}.
\]
Резолюция с основното правило за $predok$ дава
\[
\{\neg roditel(maria,petar)\}.
\]
Този литерал се резолвира с факта $roditel(maria,petar)$ и се получава $\Box$. Следователно целта е следствие на програмата.
\end{example}

\begin{remark}
В Пролог обикновено се използва SLD-резолюция: на всяка стъпка се избира подцел, унифицира се с главата на правило или факт и се заменя с тялото на правилото. Натрупаните унификатори дават отговора за променливите в заявката.
\end{remark}

\section{Изпитен фокус}

\begin{itemize}
\item Да се дефинират: литерал, елементарна дизюнкция, КНФ, дизюнкт, верен дизюнкт в структура при оценка, тъждествено верен дизюнкт и изпълнимо множество от дизюнкти.
\item Да се знае, че $\Box$ е празният дизюнкт и е неверен при всяка интерпретация.
\item Да се дефинира непосредствена резолвента:
\[
R_L(D_1,D_2)=
(D_1\setminus\{L\})\cup(D_2\setminus\{\overline L\}).
\]
\item Да се дефинират субституция, унификатор и най-общ унификатор.
\item Да се дефинира предикатна либерална резолвента чрез преименуване на променливите, унификация на противоположни атоми и прилагане на унификатора към остатъка от дизюнктите.
\item Да се дефинира резолютивен извод и означението $\Gamma\vdash_r\delta$.
\item Да се формулират теоремите за коректност и пълнота:
\[
\Gamma\vdash_r\Box\Rightarrow\Gamma\text{ е неизпълнимо},
\qquad
\Gamma\text{ е неизпълнимо}\Rightarrow\Gamma\vdash_r\Box.
\]
Пълнотата в този вид се отнася за първоредов език без формално равенство; при равенство са нужни допълнителни аксиоми или правила.
\item Да се възпроизведе доказателството за коректност: индукция по стъпките на извода, използваща резолютивната лема; празният дизюнкт не може да бъде удовлетворен.
\item Да се обясни ролята на Ербрановите частни случаи и унификацията при предикатната резолюция.
\item Да се решава задача за превеждане на предикатни формули в дизюнкти и за извеждане на $\Box$.
\item Да се разпознават факти, правила и цели като хорнови дизюнкти и да се обясни връзката им с Пролог и SLD-резолюцията.
\end{itemize}
